\documentclass[11pt]{article}
\usepackage[a4paper, margin=1cm]{geometry}
\usepackage[usenames,dvipsnames,svgnames,table]{xcolor}
\usepackage{graphicx}
\usepackage{fancyhdr}
\usepackage{background}
\usepackage{xltxtra}
\usepackage{fontspec}
\setmainfont[Mapping=tex-text]{Linux Libertine O}
%\usepackage[utf8]{inputenc}
\usepackage[frenchb]{babel}


\pagestyle{empty}
\thispagestyle{empty}

\title{\textcolor{MidnightBlue}{\textbf{AGGRASTAT dosage et  administration}}}
\author{\textcolor{MidnightBlue}{\textsc{Yahyaoui} Mohamed Kaddour, \textsc{MD}}\\
		\textcolor{MidnightBlue}{Service de cardiologie, EH Dr.Benzerdjeb}}
\date{\textcolor{MidnightBlue}{\today}}
\begin{document}

\maketitle

\section{Administration}

Mettre les 50cc du flacon dans 250cc de SGI ou SSI.\\
Bien mélanger \\
\subsection{Pas d'angioplastie prévue avant 04 heures}
\begin{itemize}
\item{dose de charge: 0.4 microgram/kg/min pendant 30 minutes.}
\item{dose d'entretien: 0.1microgram/kg/min les 48heures restantes.}
\end{itemize}
\subsection{Syndrome coronaire aigu}
Dans ce cas là l'Aggrastat doit être administré conjointement avec \emph{l'héparine non fractionnée!} \\
\begin{itemize}
\item{dose de charge: 25microgram/kg sur 3 minutes. \hfill initiée avec le début de l'angioplastie.
\item{dose d'entretien: 0.15 microgram/kg/min pendant 18-24h voir 48h.} }
\end{itemize}
\section{Précautions}
\begin{itemize}
\item{les patients les plus à risque de saignement sont :
    \begin{itemize}
        \item{faible masse corporelle}
        \item{les femmes}
        \item{sujets âgés}
    \end{itemize}}
\item{Dans ces cas là, surveiller l'effet de l'héparinothérapie}
\end{itemize}
\subsection{incompatibilités}
l'Aggrastat est incompatible avec le \textbf{Diazepam}, \hfill ne pas l'administrer dans la même voie d'abord.

\subsection{Insuffisance rénale}
en cas de Cl\MakeLowercase{créatinine} < 30ml/min : diminuer les doses d'Aggrastat de \textbf{50\%}.
\section{Surdosage}
en cas de surdosage:
\begin{itemize}
\item{arrêter immédiatement la perfusion dès les premiers signes de saignement, le plus souvent oral.}
\item{transfusion de sang si celà s'avére nécessaire.}
\item{l'Aggrastat peut être épuré, le cas échéant par \textbf{hémodialyse}}
\end{itemize}

\backgroundsetup{
  scale=1,
  color=black,
  opacity=1,
  angle=0,
  position=current page.south,
  vshift=60pt,
  contents={%
  \small\sffamily%
  \begin{minipage}{.8\textwidth}
  \parbox[b]{.6\textwidth}{%
    Cardiologie}\hfill
  \parbox[b]{.4\textwidth}{%
    \raggedleft EH Dr.Benzerdjeb,Aïn Témouchent}\\
  \textcolor{MidnightBlue}{\rule{\textwidth}{1.5pt}}\\
  \end{minipage}\hspace{.02\textwidth}%
  \begin{minipage}{.18\textwidth}
  \includegraphics[width=\linewidth,height=15pt,keepaspectratio]{logoEH1}
  \end{minipage}%
  }
}

\thispagestyle{empty}
\end{document}